% THETA3-B23 contribution by /root/definite_bridge, 21 November 2026.
\subsection{THETA3-B23: an algebraic inverse on a full CM Hodge block}\label{B23}

The finite evaluation inverse in A5 has a motivic replacement on
a full CM Hodge block.  The replacement is an algebraic
correspondence multiplied by one explicitly identified positive
comparison scalar.  It is common to every vector in that block.
This statement does not require the Hodge conjecture for the ball
surface, and does not assert that the image of one fixed B17
kernel is a rational Hodge block.

\paragraph{The CM source and the geometric normalization.}
Let $X/\overline\Q$ be the smooth projective ball surface, or the
finite disjoint union of projective PEL components at the level
used in A5.  A further finite principal cover is allowed, with
its degree retained in pullback and trace.  Let $A/\overline\Q$
be a CM abelian variety and let
\[
 T\subset H^2(A,\Q),\qquad
 \alpha:T\longrightarrow H^2(X,\Q)
 \tag{B23a}
\]
be an algebraically projected rational Hodge summand and an
algebraic correspondence.  Assume that $T$ is a sum of non-Tate
simple Hodge structures.  For the prescribed finite family,
032, Corollary 6.7, printed p.45, gives finitely many rational
Hodge maps on degree-one cohomology from the specified CM sources.
Realize them by morphisms
$X\to B_i$ and combine these into $f:X\to A=\prod_iB_i$.
Then $\alpha=f^*$ on the non-Tate part of $H^2(A)$ contains
all the prescribed theta wedges in its $(2,0)$ image.
This is a full-block enlargement inside their actual CM cup
source; it does not require that the entire Albanese be CM.
Equivalently one may use the cup image of the CM part of
$H^1(\operatorname{Alb}X,\Q)$.  The degree-one realization
is 032, Lemma 2.2, printed p.5, and its use for the CM source
maps is in Section 6.  We use the accepted CM
Hodge conjecture for products of CM abelian varieties.

Put $M=\operatorname{im}\alpha$.  On holomorphic two-forms use
the geometric Hermitian pairing
\[
                  (f,g)_{\rm geom}=\int_X\overline f\wedge g.
 \tag{B23b}
\]
It is positive with the complex orientation.  The Haar
normalization in a particular automorphic realization is
recorded by a fixed positive constant $\nu_G$:
$\langle f,g\rangle_G=\nu_G(f,g)_{\rm geom}$.
One may choose geometric normalization, when $\nu_G=1$.
Finite covering degrees or fixed rational component weights
can instead be included in (B23b).  No unknown factor is
absorbed into a character-dependent normalization.

We retain A5's full adelic $\U(V)$ realization, its fixed
splitting, central transformation law, finite level, and scalar
archimedean canonical type.  If starting with a form on one
principal $\SU(V)$ component, first use the deeper principal
level and component realization of B12 and 032, Lemma 5.2,
printed pp.27--28; extend by zero on the other components.
All CM subspaces and projectors below refer to this realized
finite collection of algebraic components.  Thus the restriction
operators of A5 apply to $\mathcal E_\chi$; an unspecified
extension of an arbitrary $\SU(V)$ form is not being assumed.

\paragraph{An algebraic retract of the CM source.}
Every non-Tate simple summand of $M$ is primitive.  Indeed,
cup product with an ample class is a Hodge map
$M(-1)\to H^4(X,\Q)\simeq\Q(-2)$ on each component,
or equivalently $M\to\Q(-1)$ after untwisting.  A
non-Tate simple Hodge structure has no such quotient.
Consequently $q_X(x,y)=-\int_Xx\cup y$ is a polarization
on $M$.  Choose a rational polarization $q_T$ on $T$.
Define the polarization adjoint $\alpha^\dagger$ using $q_T$
and the nondegenerate middle cup pairing on $H^2(X)$.
This adjoint is algebraic: the transpose of the given
correspondence is algebraic, and the inverse polarization
tensor on the CM summand $T$ is a Hodge tensor on a product
of $A$, hence algebraic by the accepted CM Hodge conjecture.
Only this CM-source inverse is needed; there is no inverse
Lefschetz or Hodge-conjecture assumption on $X\times X$.

For clarity, all cycles needed here may be defined over
$\overline\Q$.  Spread a cycle over $\C$ representing a
specified rational Hodge tensor over a finite-type
$\overline\Q$ parameter scheme, restrict to a smooth connected
open on which its components are flat, and specialize at a
$\overline\Q$ point.  On the constant ambient abelian variety
the cycle class is locally constant in that connected family.
The specialization therefore has the same class.  The usual
algebraic degree projectors of an abelian variety are also
available, for example by polynomial interpolation in
multiplication by an integer greater than one.

Set $S=\alpha^\dagger\alpha$.  The positive Hermitian forms
obtained from the two polarizations make $S$ self-adjoint and
positive semidefinite, with
$\ker S=\ker\alpha$.  Thus it is semisimple.  If $r(t)\in\Q[t]$
is the product of the nonzero distinct factors of its minimal
polynomial, then $r(0)\ne0$.  Choose $a,b\in\Q[t]$ satisfying
$ta(t)+r(t)b(t)=1$ and put
\[
 d(t)=a(t)-\frac{a(0)}{r(0)}r(t),\qquad
 \beta=d(S)\alpha^\dagger,
 \qquad e_M=\alpha d(S)\alpha^\dagger.
 \tag{B23c}
\]
The zero image case is harmless and is discarded.  The
polynomial $d$ vanishes at zero and is $t^{-1}$ on the
nonzero spectrum of $S$.  Hence $e_M$ is the orthogonal
cohomological projector onto $M$, and $\beta$ is a section
of $\alpha$ on $M$.  Every operator in (B23c) is algebraic.
This is a statement about the action of correspondences on
cohomology; it makes no assertion of an idempotent identity
modulo rational equivalence in a Chow ring.

Moreover, every rational Hodge endomorphism $u$ of $M$ is
algebraic.  The map $\widetilde u=\beta u\alpha$ is a Hodge
endomorphism of the CM summand $T$, so is algebraic by the
same CM Hodge assumption.  The correspondence
$\alpha\widetilde u\beta$ acts as $u$ on $M$ and as zero
on its orthogonal complement.  This argument also works
with algebraic coefficients.  In particular, all the
isotypic projectors and matrix units used below are
realized by algebraic correspondences on $X$.

\paragraph{One simple type with its full multiplicity.}
Fix a non-Tate simple CM Hodge structure $S_0$ occurring in
$M$, with multiplicity $m$, and write its isotypic component
as $S_0^m$.  Its endomorphism field $K$ is a CM field and
$S_0$ has rank one over $K$.  Here is the elementary
representation-theoretic reason.  A simple rational
representation of the commutative Mumford--Tate group has
commutative division endomorphism algebra, and is
one-dimensional over that field.  The positive involution
from a polarization is complex conjugation.  If it were
trivial, the Hodge circle would have image in a compact
subgroup of a product of real multiplicative groups and
would be trivial.  A weight-two simple structure would
then be Tate.  This excludes the totally real case here.

Choose an embedding $\chi:K\hookrightarrow\overline\Q$
whose eigenspace has Hodge type $(2,0)$, and put
\[
 \mathcal E_\chi=(S_0^m\otimes_\Q\C)_\chi
                  \subset H^0(X,K_X),
 \qquad \dim_\C\mathcal E_\chi=m.
 \tag{B23d}
\]
The projector $e_\chi$ is allowed algebraic coefficients.
It is algebraic by (B23c), since
$\operatorname{End}_{\rm Hdg}(S_0^m)=M_m(K)$ and
$K\otimes_\Q\overline\Q$ splits into its embedding factors.
Distinct isotypic components and distinct embedding
components are orthogonal for the Hermitian polarization.
For the latter assertion, multiplication by $a\in K$ has
adjoint multiplication by $\overline a$; choosing $a$ whose
two embedding values differ forces the cross pairing to
vanish.

Choose compatible Betti and algebraic de Rham bases using
one copy of $S_0$ and the algebraic isomorphisms to its $m$
copies.  If $\beta_1,\ldots,\beta_m$ are the holomorphic
representatives of the resulting algebraic-coefficient
Betti classes and $\omega_1,\ldots,\omega_m$ the de Rham
basis, compatibility with comparison gives one scalar
$P_\chi\in\C^\times$ such that
\[
 \omega_j^{\rm an}=P_\chi\beta_j\quad(1\leq j\leq m),
 \qquad
 G_\omega=\nu_G|P_\chi|^2G_B,
 \quad
 (G_B)_{ij}=\int_X\overline{\beta_i}\wedge\beta_j.
 \tag{B23e}
\]
The matrix $G_B$ is positive definite with entries in
$\overline\Q$: cup product is rational on Betti cohomology,
and the classes and their conjugates have algebraic
coefficients.  The $\omega_j$ are actual algebraic
canonical forms.  On a smooth projective surface the
$F^2$ piece of algebraic de Rham cohomology is
$H^0(X,K_X)$, and the algebraic projector respects this
filtration.  The common period in (B23e) is a consequence
of the copy isomorphisms, not of independent rescaling of
the individual forms.

\paragraph{The finite evaluation rows and their actual fibre frames.}
Fix a rational negative line $L\subset V$ and its rational
orthogonal basis as in B12.  Use A5 for the full block
$\mathcal E_\chi$.  List all finite definite double cosets
occurring in its finitely many chosen translates, retaining
repetitions when they are different rows.  Their evaluation
points have the form $x_\ell=(L,b_\ell)$ with finite
adelic $b_\ell$.  These points and the finite Hecke branches
joining them are defined over $\overline\Q$.

There are direct proofs in this particular PEL model.
At the rational line, the three rational orthogonal
projectors preserve the Hodge filtration.  Adding integral
multiples of their graphs to the finite-type moduli
problem makes the chosen point isolated: a negative line
preserved by the projectors must be the specified negative
coordinate line.  An isolated complex point of a
finite-type scheme over a number field is algebraic.
This is precisely the argument in 032, Lemma 5.3,
printed pp.29--30.  For a finite adelic modification of
the lattice, clear denominators by an integer.  The
corresponding isogeny has a finite, $\mathcal O_E$-stable
torsion kernel.  Such kernels and finite level structures
are parametrized by finite etale schemes in characteristic
zero; their quotients and polarization homomorphisms are
algebraic.  The target lattice has the prescribed
polarization type because the adelic modification is
unitary.  Both projections of the resulting Hecke
correspondence become finite etale after passing to a
common level.  The quotient and polarization descent are
the same constructions written explicitly for the split
spherical correspondences in 032, Section 5.3,
printed pp.30--32.  This argument allows the finitely many
adelic translates used by A5, not just one split prime.

Let $\kappa_\ell$ denote the analytic canonical-fibre frame
used to define the automorphic scalar at $x_\ell$.  Choose
one algebraic frame $\eta_0$ at the base point and transport
it through the algebraic differentials of these finite
Hecke branches.  Since the two maps on the marked ball
chart have the same archimedean coordinate, their
differentials also identify the frames $\kappa_\ell$.
Consequently the transported algebraic frames satisfy
\[
                  \eta_\ell=t_x\kappa_\ell
                  \quad\hbox{with one }t_x\in\C^\times.
 \tag{B23f}
\]
If independent algebraic choices of the $\eta_\ell$ are
preferred, their nonzero algebraic ratios are retained
in the row matrix below.  Formula (B23f) does not assert
that a Betti ball-coordinate frame is algebraic.

With probability Haar measure on each definite quotient,
let $D$ be the diagonal matrix of the finite double-coset
masses.  Its positive entries are rational: they are
reciprocal finite stabilizer orders divided by their
finite sum.  Rows forced to vanish by a stabilizer
character can be removed.  Define
\[
 A_{\ell j}=\frac{\omega_j(x_\ell)}{\eta_\ell}
       \in\overline\Q,
 \qquad R=t_x A,
 \qquad H=A^*DA\in M_m(\overline\Q).
 \tag{B23g}
\]
The finite family in A5 separates $\mathcal E_\chi$, so
$A$ has full column rank and $H$ is positive definite.

\paragraph{The canonical-fibre period, including the Tate factors.}
The number $t_x$ has a specified CM-period meaning.  Use
the actual reference PEL abelian scheme of 032, Section
5.1.  Its active Hodge bundles have ranks two and one;
write them as $\mathcal P=\omega_\mu$ and
$\mathcal L=\omega_{c\mu}$.  Gauss--Manin and the
polarization give algebraic isomorphisms
\[
 \operatorname{KS}:\mathcal P\otimes\mathcal L
                         \simeq\Omega_X^1,
 \qquad
 \det\mathcal P\otimes\mathcal L^{\otimes2}
                         \simeq K_X.
 \tag{B23h}
\]
The other embedding summands are constant along the ball
and contribute no tangent directions.  To check that
the active map is an isomorphism, differentiate the
negative line and its annihilator in a marked period
chart.  The two independent coordinates of the line
give the two independent components of
$\mathcal P\otimes\mathcal L$ under the nondegenerate
polarization pairing.  Thus its determinant is nonzero.
This also proves (B23h) for the higher-degree CM field,
without importing a dimension-three assertion globally.

At $x$, choose algebraic bases of $\mathcal P_x$ and
$\mathcal L_x$ and write their comparisons with the
rational graph bases as a matrix $\Omega_\mu$ and a
scalar $P_{c\mu}$, respectively.  The graph bases and
the derivative of the rational line have algebraic
entries.  The de Rham polarization on $H^1$ compares
with its Betti alternating pairing by $(2\pi i)^{-1}$.
There are two such pairings in the determinant of
Kodaira--Spencer.  Hence, for the frame obtained from
these algebraic bases,
\[
 t_x=\zeta_x\,
       \frac{\det(\Omega_\mu)P_{c\mu}^{\,2}}
                         {(2\pi i)^2},
                 \qquad\zeta_x\in\overline\Q^\times.
 \tag{B23i}
\]
Here $\zeta_x$ is the explicitly determined determinant
of the algebraic graph, polarization, and tangent-frame
matrices.  It can be included in $A$ instead if desired.
The factor $(2\pi i)^{-1}$ can also be checked by pairing
two de Rham one-classes with the $(g-1)$st power of an
algebraic polarization class: its comparison contributes
$(2\pi i)^{g-1}$, while the top-degree de Rham trace
contributes $(2\pi i)^{-g}$.

An exact primary normalization check is de Shalit--Goren,
\emph{A theta operator on Picard modular forms modulo
an inert prime}, Section 1.8, Proposition 1.22 and
Section 1.11.2, Lemma 1.27 and Proposition 1.28:
equation (1.141) is
$\operatorname{KS}(d\zeta_1\otimes d\zeta_3)
=-\delta(2\pi i)^{-1}du$ and
$\operatorname{KS}(d\zeta_2\otimes d\zeta_3)
=(2\pi i)^{-1}dz$.
The complete displayed computation and its algebraic
interpretation in Section 1.11.3 give precisely the
power in (B23i).  The version on the authors' site has
these statements on printed pp.26--27 and 31--33;
the published article uses the same numbered statements.
No identification $\det\mathcal P\simeq\mathcal L$
with an unspecified scalar is used here.

\paragraph{Algebraic matrix and common positive scalar.}
In the coherent basis $\omega$, the operator of A5 and
its inverse are exactly
\[
 \begin{split}
 Q_\chi&=
     \frac{|t_x|^2}{\nu_G|P_\chi|^2}\,G_B^{-1}H,\\
 Q_\chi^{-1}&=s_\chi K_\chi,
 \qquad
 s_\chi=\frac{\nu_G|P_\chi|^2}{|t_x|^2}>0,
 \qquad K_\chi=H^{-1}G_B\in M_m(\overline\Q).
 \end{split}
 \tag{B23j}
\]
Indeed $R^*DR=|t_x|^2H$ and (B23e) is the domain
Gram matrix; the definition of the adjoint gives the
first equality and inversion gives the second.
Because this is the full multiplicity of a single
embedding component, the matrix $K_\chi$ extends by
zero to an element of
$\operatorname{End}_{\rm Hdg}(M)\otimes_\Q\overline\Q$.
The CM-source retract above therefore realizes it as
an algebraic correspondence on $X$.  On the selected
complex block it is positive and self-adjoint for
the geometric pairing, since it is a positive real
multiple of $Q_\chi^{-1}$.  Positivity of all algebraic
conjugates is neither asserted nor needed.

For every $f,g\in\mathcal E_\chi$, A5 now reads
\[
 \langle f,g\rangle_G
   =s_\chi\sum_{j,\alpha}
       \overline{\mathcal P_{j,\alpha}(K_\chi f)}
                       \mathcal P_{j,\alpha}(g).
 \tag{B23k}
\]
The periods are the actual totally definite $\U(2)$
Bessel periods with A5's fixed translates and probability
Haar measures.  The same $K_\chi$ and $s_\chi$ work for
every source character whose image lies in this block.
When two selected CM tensor lines have the same complete
Hodge type, the algebraic CM comparison of A10 places
their nonzero images in this same isotypic block and
embedding component.  This is the appropriate meaning
of a common normalization for balanced pairs.

There is also an exact one-reference calibration.
Choose $0\ne f_0\in\mathcal E_\chi$ once and put
$B(f,g)=\sum_{j,\alpha}
\overline{\mathcal P_{j,\alpha}(K_\chi f)}
\mathcal P_{j,\alpha}(g)$.  Then $B(f_0,f_0)>0$ and
\[
 \langle f,g\rangle_G
       =\frac{\langle f_0,f_0\rangle_G}{B(f_0,f_0)}
                           B(f,g).
 \tag{B23l}
\]
If $f_0$ is specified by an algebraic-coefficient Betti
class its geometric squared norm is algebraic.  The
denominator in (B23l) still contains the common period;
this formula does not claim that $s_\chi$ is algebraic.
Formula (B23i) identifies the exact remaining
Kodaira--Spencer CM tensor, rather than hiding it in
an arbitrary evaluation normalization.

\paragraph{What block projection repairs, and what it leaves.}
Evaluation at a point is not a morphism of weight-two
Hodge structures.  Before projection, the finite
evaluation Gram matrix can have nonzero cross terms
between distinct CM blocks, although the geometric
cup pairing has none.  Replace each restriction by
$R_ae_\chi$ and apply (B23j) separately on its full
block.  The algebraic projectors remove these cross
terms exactly.  Averaging over a CM action alone
would only force a commuting operator; even in
multiplicity one its positive scalar remains an
analytic norm ratio.  No algebraicity of that scalar
follows from the commutation relation.

The assertion proved here concerns the full CM block
with its algebraic global forms and all selected
finite evaluation rows.  A particular image
$\mathcal E_0$ of the B17 integral operator need not
be a rational Hodge summand, and its adjoint involves
an integral over the original source group.  Neither
algebraicity of that integral Gram matrix nor equality
of its inverse with $K_\chi$ has been established by
the finite-frame argument.  Thus (B23k) closes the
algebraic-frame correction for A5 on a full CM block,
up to the single specified positive factor, while
the literal identification with one fixed arithmetic
\texttt{MainX}/\texttt{WeilTensorModels} package remains
a separate kernel and normalization comparison.
In particular, the coordinates of an original unnormalized
theta vector in the coherent basis $\omega$ can contain
the type-dependent constants retained in B19--B21.
The algebraicity of $K_\chi$ does not make those
coordinates algebraic.

\paragraph{The Grothendieck framework.}
The supplied manuscript locators are
Grothendieck's \emph{cote 44}, scans 7--8, in the letter
to Serre dated 12 August 1964: the polarization is
constructed through a transposed cycle correspondence,
and its kernel is removed before passing to the abelian
quotient.  The inverse-correspondence and critical-degree
discussion is \emph{cote 16}, scans 43--46, especially
the equivalences (iv) and (iv bis) on scan 44 and
their transposition argument on scan 45.  Here (B23c)
constructs the particular inverse on a CM source by a
rational polynomial and the accepted CM Hodge theorem.
It therefore proves the needed homological retract;
it does not invoke the general standard conjectures
or replace the distinction between a cohomological
projector and a Chow idempotent.
